\section{Using the generator to evaluate post-hoc explanations}
\label{sec:generator-post-hoc-explainability}

The generator makes it possible to evaluate whether a forecast explanation
highlights the inputs that appear in a known data-generating equation. We
specify the equations and coefficients before generating any trajectories, so
each scenario has an explicit reference: the set of direct parents of the
target. A parent is a lag in the univariate case and a channel--lag pair in
the multivariate case. The notebooks under
\texttt{experiments/chronos2\_explainability/} display the equations and record
the seeds, generator settings, and model revision.

\subsection{Forecast and evaluation protocol}

For each equation or system, we generate eight trajectories with independent
Gaussian innovations and evaluate six windows per trajectory. Each window
contains 128 observations. We use the frozen Amazon Chronos-2 checkpoint
\citep{ansari2025chronos2} to forecast the one-step-ahead median for one target
channel. We compare the gradient of that forecast with respect to each
observed input and the same gradient multiplied by the input value. We compute
both on the original input scale. In each scenario, we check the differentiable
forecast against the official pipeline and check selected derivatives by
finite differences.

If the equation has $k$ direct parents, we rank the inputs by absolute
attribution and count how many of those parents appear in the top $k$.
The count is $k$ times recall at $k$. Ties at the cutoff receive the expected
number of hits under uniform tie breaking. We average the count over windows
within each trajectory and then over trajectories. In the multivariate case,
an attribution must identify both the correct channel and the correct lag.

\subsection{Recovery of direct dependencies}

The first two notebooks use fixed linear equations chosen to keep propagated
effects below the direct effects. Table~\ref{tab:post-hoc-top-k} reports their
average counts. Plain gradients recovered more direct parents in all four
scenarios. Even in the best scenario, the average was about three of seven.
These measurements describe the agreement between the Chronos-2 explanations
and the specified generator equations. An explanation can reflect the
forecast's sensitivity accurately while disagreeing with the generator.

\begin{table}[htbp]
\centering
\caption{Mean number of direct parents recovered in the top $k$.}
\label{tab:post-hoc-top-k}
\begin{tabular}{lrrr}
\toprule
Scenario & $k$ & Gradient & Gradient $\times$ input \\
\midrule
Univariate, consecutive lags & 7 & 3.0625 & 2.1250 \\
Univariate, sparse lags & 7 & 1.8542 & 1.2292 \\
Multivariate, cross-channel parents & 2 & 0.1250 & 0.0208 \\
Multivariate, own and cross-channel parents & 2 & 0.3958 & 0.3125 \\
\bottomrule
\end{tabular}
\end{table}

\subsection{What the direct reference leaves out}

The third notebook varies mediation in a three-channel system,
\begin{align*}
A_t &= \varepsilon_t^A, \\
B_t &= b A_{t-1} + \varepsilon_t^B, \\
Y_t &= 0.8 B_{t-1} + d A_{t-2} + \varepsilon_t^Y.
\end{align*}
Increasing $A_{t-2}$ by one yields a controlled direct effect of $d$ when
$B_{t-1}$ is held fixed. With the innovations held fixed and the mediator
allowed to respond, the effect on $Y_t$ is $d+0.8b$. The experiment uses
$(b,d)=(0.05,0.05)$, $(0.8,0.05)$, and $(0.8,0)$. In the final case,
$A_{t-2}$ has an effect of $0.64$ through $B_{t-1}$ and receives no credit
under the direct-parent reference, even if an explanation highlights it.

This example explains why the direct-recovery experiment selects equations
with controlled propagation. Its score answers which inputs enter the target
equation directly. It cannot establish whether an explanation recovers
influence transmitted through earlier states. A complementary evaluation could
derive a reference for propagated effects by intervening in the known
generator and allowing descendants to respond. The intervention and forecast
horizon would then define precisely which influence the score measures.
